% Avsnitt 17.10, ur lectures/L17/appendix/b_exercises.md.

\section{Övningar}
\label{c17:sec:exercises}

\begin{aside}
När din mottagningslogik är klar verifierar \file{controller/can_controller_tb.vhd} sändning och
mottagning tillsammans, och \code{make build-project} slutar hoppa över den: alla sex testbänkar i
\file{controller/} körs härifrån och framåt. De två i \file{bridge/} fortsätter hoppas över tills
\chapref{c18:ch} och \ref{c19:ch} skriver modulerna de prövar.

Om den fortfarande rapporteras som överhoppad efter att du tror att mottagningsvägen är klar, läs
meddelandet före din design. Bygget avgör saken genom att leta efter \chapref{c17:ch}:s CRC-grindning
på mottagningssidan inuti \file{can_controller.vhd}, så en korrekt kontroller som stavar de två
uttrycken annorlunda hoppas över i stället för att köras. \code{CI_BUILD_ALL=1 make build-project}
tvingar fram den, och bilaga~\ref{app:sim} förklarar varför kontrollen fungerar så.
\end{aside}

\exercise{Resonemang om tajmingfixen på RX-sidan}{Resonemang}
\label{c17:ex:timing}
\xpart{a} \Secref{c17:sec:preempt} visar att \code{rx_shift_reg} måste aktiveras
\emph{kombinatoriskt}, under själva \code{STATE_LOAD_*}-tillståndet, i stället för att vänta tills
\code{state} synligt blir skifttillståndet. Förklara med egna ord varför det att först vänta in att
tillståndsövergången blir klar skulle få \code{rx_shift_reg} att missa just det sampel som utlöste
väntan från början.

\xpart{b} \code{STATE_CRC_LO_WAIT} väntar på \code{bt_bit_done}, inte på \code{bt_sample}, till
skillnad från varje väntan vid en gruppgräns i designen (tillstånden \code{STATE_LOAD_*}). Förklara
med egna ord varför just den övergången måste ligga i linje med en \emph{bitperiodsgräns} och inte
med en sampelpunkt. Vilka andra tillstånd delar utlösaren \code{bt_bit_done}, och vad har de
gemensamt med den?

\xpart{c} Anta (hypotetiskt) att \code{STATE_CRC_WAIT} också väntade på \code{bt_sample} för
RX-rollen, utöver att \code{STATE_LOAD_CRC_HI} gör detsamma omedelbart efteråt. Förklara vad som
skulle gå fel, och varför det inte så mycket skulle visa sig som ett rent haveri som att en mottagande
nod hamnar en bitperiod längre efter en sändande nod än den borde.

\exercise{Arbitrering för hand}{Resonemang}
\label{c17:ex:arbhand}
Två noder begär sändning på samma cykel: nod A med ID \code{0x2A0}, nod B med ID \code{0x2C0}.

\xpart{a} Skriv båda ID:na binärt (11 bitar, MSB först) och markera den första bitposition där de
skiljer sig.

\xpart{b} Säg med \secref{c10:sec:arbitration}:s arbitreringsregel (dominant \code{0} vinner) vilken
nod som vinner och vid vilken bit förloraren faller ur.

\xpart{c} Den förlorande noden upptäcker sin förlust med
\code{bt_sample = '1' and txsr_tx_bit = '1' and rx_bus_s2 = '0'}. Säg, för förloraren vid den biten,
vad var och en av de tre termerna evalueras till och varför villkoret slår till.

\xpart{d} Förloraren sätter \code{error <= '1'} och faller ner till \code{STATE_IDLE} medan vinnaren
fortfarande driver bussen. \code{error} nollställs i \code{STATE_START}. Förklara varför förlorarens
\code{error} skulle bli oobserverbart för en anropare om \code{STATE_IDLE} gick in i en mottagning på
en dominant \emph{nivå}.

\xpart{e} En lockande lagning är att i stället leta efter \emph{flanken} från recessiv till dominant,
i linje med hur \secref{c10:sec:frameformat} beskriver SOF. Visa att det egentligen inte fungerar,
med hjälp av stoppbitsregeln: vad sätter sändaren in efter fem dominanta bitar i rad, och vad gör
bussen på biten efter det? Ungefär hur långt in i vinnarens ram hinner förloraren innan flanktestet
slår till ändå?

\xpart{f} Ange regeln designen faktiskt använder (\secref{c16:sec:roles}). Två fakta gör den
hållbar; ge båda: vad stoppbitsregeln begränsar inne i det stoppade området, och vad den ostoppade
svansen lägger till ovanpå det som skjuter tröskeln till elva. Säg sedan varför tomgångsräknaren
måste nollställas till \textbf{fullt} och inte till tomt när noden kommer ur reset.

\exercise{SOF-biten som ingen behåller}{Resonemang}
\label{c17:ex:sof}
\code{STATE_SOF} aktiverar \code{rx_shift_reg} under en bit och kastar sedan värdet.
\Secref{c17:sec:states} förklarar varför; den här övningen handlar om att se felet med egna ögon.

\xpart{a} Ta ID \code{0x000}, DLC \code{0}. Skriv ut bitarna en sändare lägger på ledningen från SOF
till och med kontrollfältets slut, med \secref{c16:sec:fields}:s fälttabell, och markera var
\code{tx_shift_reg} sätter in varje stoppbit. Kom ihåg att följdräknaren räknar med SOF.

\xpart{b} Gör nu samma sak sett från en mottagare som hoppade över SOF, så att dess följdräknare
ligger en dominant bit efter. Markera var \emph{den} förväntar sig stoppbitar. Ställ upp de två mot
varandra och ange den första bitposition där de är oense.

\xpart{c} Vid den positionen skickar sändaren en stoppbit och mottagaren väntar sig ingen. Vad gör
mottagaren med den, och vad gör det av varje bit efter den?

\xpart{d} Upprepa \textbf{a)} och \textbf{b)} för ID \code{0x123}, det ID som
\code{can_controller_tb} använder i fall 1, och bekräfta att de två räknarna är överens hela vägen.
Ange den allmänna regeln för vilka identifierare som går sönder: vad måste gälla för de första
identifierarbitarna, och hur många av de 2048 standard-ID:na uppfyller det?

\xpart{e} Fall 2 i \code{can_controller_tb} skickar \emph{faktiskt} ID \code{0x000}. Förklara varför
buggen ändå inte skulle visa sig där, och säg vad du skulle lägga till i testbänken för att fånga
den.

\exercise{Att ta emot en ram utan data}{Resonemang}
\label{c17:ex:dlc0}
Gå igenom \textbf{mottagningsrollen} genom en komplett ram med \code{DLC = 0}, med hjälp av tabellen
per tillstånd i \secref{c17:sec:states}. En mottagare vet inte DLC:n förrän den har avkodat
kontrollfältet, så det här är också den väg som prövar om din övergång i \code{STATE_CTRL} är rätt.

\xpart{a} Lista varje tillstånd noden passerar, från \code{STATE_IDLE} tillbaka till
\code{STATE_IDLE}. Markera den enda övergång som skiljer sig från en ram som bär data, och säg vad
som får den att slå till.

\xpart{b} Ange för varje \code{STATE_LOAD_*}-tillstånd i din lista det \code{rxsr_bit_count} som det
kombinatoriska nätet föregriper, och säg vilket \code{bt_sample} det föregripandet måste vara på
plats för: det som utlöser övergången ut ur \code{STATE_LOAD_*}, eller nästa.

\xpart{c} Anta att föregripandet vore en cykel för sent överallt, alltså att \code{rxsr_enable} gick
hög först när \code{state} synligt läste skifttillståndet. Hur många riktiga bitar skulle
\code{rx_shift_reg} ha missat för en ram med \code{DLC = 0} när noden når \code{STATE_CRC_LO_WAIT},
och vad skulle \code{crc_valid} läsa där?

\xpart{d} Ditt svar på \textbf{c)} är ett CRC-fel som rapporteras många bitperioder efter det
faktiska misstaget. Relatera det till \secref{c17:sec:preempt}:s påstående att en felställning på en
cykel ackumuleras till en tyst felställning som visar sig först senare, långt från sin orsak, och säg
varför det gör den här sortens bugg dyr att felsöka enbart utifrån testbänkens utskrifter.
